\documentclass[11pt]{article}

\usepackage[margin=1in]{geometry}
\usepackage{amsmath,amssymb,amsfonts}
\usepackage{amsthm}
\usepackage{bm}
\usepackage{mathrsfs}
\usepackage{hyperref}
\usepackage{enumitem}
\usepackage{listings}
\usepackage{xcolor}
\usepackage{tikz}

\definecolor{codegray}{rgb}{0.15,0.15,0.15}
\definecolor{codebg}{rgb}{0.97,0.97,0.97}
\definecolor{codekw}{rgb}{0.10,0.10,0.60}
\definecolor{codestr}{rgb}{0.50,0.10,0.10}
\definecolor{codecm}{rgb}{0.25,0.45,0.25}

\lstdefinelanguage{TypeScript}{
  keywords={type, interface, export, import, from, const, let, async, await,
    class, implements, extends, constructor, new, return, if, else, try, catch},
  keywordstyle=\color{codekw}\bfseries,
  ndkeywords={number, string, boolean, Promise, Map, Record},
  ndkeywordstyle=\color{codekw},
  identifierstyle=\color{codegray},
  comment=[l]{//},
  commentstyle=\color{codecm}\itshape,
  stringstyle=\color{codestr},
  sensitive=true
}

\lstdefinelanguage{Python3}{
  language=Python,
  keywordstyle=\color{codekw}\bfseries,
  commentstyle=\color{codecm}\itshape,
  stringstyle=\color{codestr},
}

\lstset{
  backgroundcolor=\color{codebg},
  basicstyle=\ttfamily\small,
  breaklines=true,
  tabsize=2,
  showstringspaces=false,
  frame=single
}

\title{PIRTM, Phase Mirror, and ADR--028:\\
Lie--Algebraic Trust Weighting and Monoidal Ensemble Aggregation}
\author{Citizen Gardens \\ The Foundation of Multiplicity}
\date{April 16, 2026}

\begin{document}

\maketitle

\begin{abstract}
This report consolidates the recent developments around Phase Mirror's trust
architecture, the PIRTM operator atlas, and ADR--028's three--factor weight
model. It formalizes the Lie--algebraic structure generated by the prime--indexed
PIRTM operators, shows how this justifies additive aggregation of emergent
signals, and documents the corresponding TypeScript and Python changes in the
Phase Mirror codebase.\footnote{Trust module scaffolding and reputation system:
summary in \texttt{trust/reputation} and associated README.\,[file:12]}
It also situates these changes within Phase Mirror's three architectural planes
(PMD methodology, L0--L2 compute, and the 6--layer trust/defense system) and the
L0 invariant constraints that govern Byzantine filtering and consent.
\end{abstract}

\newpage
\tableofcontents

\section{Executive Summary}

The central development is the unification of PIRTM's multiplicity--theoretic
operator view with Phase Mirror's trust, reputation, and Byzantine filtering
mechanisms.

\begin{itemize}[leftmargin=*]
  \item The PIRTM operator atlas \(\mathcal{A} = \{U_p^{(r)}\}\) acts on a
  multiplicity Hilbert space
  \(H = \ell^2(\mathbb{P}) \otimes L^2(\mathbb{R}) \otimes \mathbb{C}^d\),
  and each bundled operator (e.g.\ \(\hat{P}_p\), \(A_{p^r}\), \(R_{p^r}\))
  is realized as a bounded operator on \(H\).\,[file:12]

  \item The set of bounded operators \(B(H)\) is an associative \(*\)--algebra
  and hence, under the commutator \([X,Y] = XY - YX\), forms a Lie algebra.
  The emergent signal
  \(\mathcal{E}_{ij}(t) = w_i^{\text{final}} w_j^{\text{final}} [\Xi_i,\Xi_j]\)
  is an element of this Lie algebra and lives in the same ambient operator
  space as its components.\,[file:12]

  \item ADR--028 declares the canonical PIRTM weight as
  \(w_i^{\text{final}} = T_i \cdot c_i \cdot r_i\), where
  \(T_i\) is an inference--tier factor, \(c_i\) is a kernel--confidence factor,
  and \(r_i\) is the existing reputation product over
  base score, stake multiplier, and consistency bonus.\,[file:12]

  \item Because \(B(H)\) is a vector space, the additive combination
  \(\mathcal{E}_{\text{ensemble}} = \sum_{(i,j)} \mathcal{E}_{ij}\) is exact
  and requires no projection back onto the atlas: closure under addition
  guarantees that the monoidal ensemble reduction is well--defined as a sum of
  commutators.\,[file:12]

  \item On the implementation side, the TypeScript reputation types and
  \texttt{calculateContributionWeight} have been extended to carry
  \emph{orthogonal} factors (\(T_i\), \(c_i\), \(r_i\)) and provenance fields,
  while preserving backwards compatibility and default semantics for existing
  callers.\,[file:12]

  \item A new L0 guard in
  \texttt{ByzantineFilter.prepareWeightedContributions} enforces the invariant
  that any Tier--3 (Moore--Penrose) cross--validated confidence requires at
  least one Tier--1 cumulant anchor in the ensemble; otherwise it is
  downgraded to a conservative floor.\,[file:12][file:6]

  \item On the Python side, a monoidal ensemble aggregator computes pairwise
  weighted commutators and reduces them via associative addition inside
  \(B(H)\), tracking composite trust and provenance for \texttt{originRecord}
  and downstream dashboards.\,[file:12][file:13]
\end{itemize}

These changes align the trust module scaffolding, the FP calibration network,
and the dissonance--oriented PMD loop with a single constitutional principle:
\emph{Phase Mirror aggregates operator provenance under constraint, not
beliefs.}\,[file:6][file:13]

\section{Phase Mirror Architectural Planes}

\subsection{Three Planes of Architecture}

The Phase Mirror system is now understood as operating across three distinct
architectural planes rather than a single conflated state space.\,[file:6]

\begin{description}[leftmargin=*,style=nextline]
  \item[Plane A --- PMD Methodology Plane.] The PMD operating loop
  \(\text{PMD} = (D, T, \pi, B, Q)\) encapsulates the dissonance--driven
  governance method: extract artifacts, map tensions, rank, produce output
  blocks, and pose a precision question.\,[file:6][file:13]

  \item[Plane B --- L0--L2 Compute Plane.] The compute tiers \(L_0, L_1, L_2\)
  partition analysis by latency and strictness: L0 foundational invariants
  (sub--100ns p99) must pass before any higher--tier work; L1 hosts rule
  evaluation and FP filtering (ms scale); L2 hosts deep semantic analysis and
  audits (tens to hundreds of ms).\,[file:6]

  \item[Plane C --- Trust/Defense Plane.] The six--layer trust and defense
  architecture (identity, economic, cryptographic, Byzantine, privacy,
  monitoring) protects the FP calibration network and the network effect
  moat, but is scoped explicitly to multi--organization operation.\,[file:6]
\end{description}

The newly introduced ADR--028 weight model and the monoidal ensemble
aggregation live at the intersection of Plane B (compute) and Plane C
(trust/defense): they are implemented at the trust module level and feed the
Byzantine calibration layer while constrained by L0 invariants.\,[file:12][file:6]

\subsection{L0 Invariants and Governance}

Phase Mirror treats L0 invariants as always--on, non--negotiable constraints
that must be validated prior to any higher--layer computation.\,[file:6][file:13]
Examples include:

\begin{itemize}[leftmargin=*]
  \item Schema hash integrity and permission bits validation for artifacts
  under analysis.\,[file:6]
  \item Nonce freshness and uniqueness for identity binding and Sybil
  resistance.\,[file:12]
  \item Drift magnitude bounds and contraction witnesses for baseline
  management.\,[file:6]
\end{itemize}

The L0 guard introduced in the Byzantine filter is a concrete instantiation of
this principle: it enforces a constitutional constraint on how Tier--3
confidence claims are permitted to enter the trust calculus.\,[file:12]

\section{PIRTM Operator Atlas and Lie Algebra Structure}

\subsection{Multiplicity Hilbert Space and Bounded Operators}

Let \(H\) denote the multiplicity Hilbert space
\[
  H = \ell^2(\mathbb{P}) \otimes L^2(\mathbb{R}) \otimes \mathbb{C}^d,
\]
where:
\begin{itemize}[leftmargin=*]
  \item \(\ell^2(\mathbb{P})\) carries prime--indexed multiplicity amplitudes.
  \item \(L^2(\mathbb{R})\) is the temporal (or spectral) component.
  \item \(\mathbb{C}^d\) is a finite--dimensional internal space (e.g.\ model
  channels or feature dimensions).\,[file:12]
\end{itemize}

The PIRTM operator atlas \(\mathcal{A}\) is a family of prime--indexed
operators
\[
  \mathcal{A} = \{ U_p^{(r)} : H \to H \mid p \in \mathbb{P},\ r \in \mathbb{N} \},
\]
including bundled forms such as \(\hat{P}_p\), \(A_{p^r}\), and \(R_{p^r}\).
Each of these is realized as a bounded operator, and hence an element of
\(B(H)\).\,[file:12]

By construction, \(B(H)\) is:
\begin{itemize}[leftmargin=*]
  \item A vector space over \(\mathbb{C}\), closed under addition and scalar
  multiplication.
  \item An associative algebra under composition \((X,Y) \mapsto XY\).
\end{itemize}

\subsection{Commutator as Lie Bracket}

Given any associative algebra \(\mathcal{A}\) over a field, defining the
commutator
\[
  [X,Y] := XY - YX
\]
yields a Lie bracket:

\begin{itemize}[leftmargin=*]
  \item Bilinearity follows from distributivity and the vector space
  structure.
  \item Antisymmetry: \([X,Y] = -[Y,X]\).
  \item Jacobi identity: \([X,[Y,Z]] + [Y,[Z,X]] + [Z,[X,Y]] = 0\),
  a standard operator identity.\,[file:12]
\end{itemize}

Therefore \(B(H)\), under the commutator, is a Lie algebra. In particular, the
set of PIRTM--generated operators \(\{\Xi_i\}\) and their commutators
\([\Xi_i, \Xi_j]\) all live in a Lie subalgebra of \(B(H)\).\,[file:12]

\subsection{Emergent Signals as Lie Algebra Elements}

For each contributing model \(i\), PIRTM produces an inferred operator
\(\Xi_i \in B(H)\) together with a final weight \(w_i^{\text{final}}\). For a
pair \((i,j)\), the emergent signal is defined as
\[
  \mathcal{E}_{ij}(t) := w_i^{\text{final}} w_j^{\text{final}} [\Xi_i, \Xi_j].
\]

Because \([\Xi_i,\Xi_j] \in B(H)\) and scalar multiplication preserves
membership in \(B(H)\), we have \(\mathcal{E}_{ij} \in B(H)\) for all \(i,j\).
Thus the emergent signal lives in the Lie algebra generated by the atlas
under the commutator, and any finite linear combination of such signals is
again in \(B(H)\).\,[file:12]

This is the key mathematical justification for treating additive combinations
of emergent signals as exact rather than approximate: no projection back onto
the atlas is required after each reduction step.

\section{ADR--028: Three--Factor Weight Model}

\subsection{Canonical Weight Decomposition}

ADR--028 defines the canonical PIRTM weight as the product
\[
  w_i^{\text{final}} = T_i \cdot c_i \cdot r_i,
\]
where:\,[file:12]

\begin{itemize}[leftmargin=*]
  \item \(T_i\) (tier factor) encodes inference provenance:
    \begin{itemize}
      \item \texttt{'cumulant-primary'} (Tier 1, exact PIRTM cumulant path),
      \item \texttt{'pseudoinverse-secondary'} (Tier 2, structured
      approximation),
      \item \texttt{'moorpenrose-tertiary'} (Tier 3, generic Moore--Penrose,
      external models only).
    \end{itemize}
  \item \(c_i\) (kernel confidence) quantifies how much of \(\Xi_i\) lies
  outside the observable subspace, via \(c_i = 1 - \|\epsilon_i\| /
  \|\Xi_i\|\) when error information is available, or via a conservative
  floor for Tier--3 when no Tier--1 anchor exists.\,[file:12]
  \item \(r_i\) (network reputation) is the existing product of base
  reputation, stake multiplier, and consistency bonus computed by the trust
  module's \texttt{ReputationEngine}.\,[file:12]
\end{itemize}

The crucial design decision is that these three factors are treated as
\emph{orthogonal dimensions} in the trust calculus: provenance (tier),
epistemic quality (kernel confidence), and social/economic behavior
(reputation) must not be collapsed prematurely into a single scalar without
retaining their separate roles and governance semantics.\,[file:12]

\subsection{TypeScript Type Extensions}

ADR--028 extends the existing reputation types with new orthogonal factors
while preserving the original fields to maintain backward compatibility:

\begin{lstlisting}[language=TypeScript,caption={ADR--028 extensions to \texttt{ContributionWeightFactors} and related types.}]
export type InferenceMethod =
  | 'cumulant-primary'
  | 'pseudoinverse-secondary'
  | 'moorpenrose-tertiary';

export type KernelConfidenceMethod =
  | 'direct'
  | 'cross-validated'
  | 'conservative-floor';

export const TIER_FACTORS: Record<InferenceMethod, number> = {
  'cumulant-primary':        1.0,
  'pseudoinverse-secondary': 0.7,
  'moorpenrose-tertiary':    0.4,
} as const;

export interface ContributionWeightFactors {
  baseReputation:   number;
  stakeMultiplier:  number;
  consistencyBonus: number;
  totalMultiplier:  number;   // r_i

  // ADR-028: orthogonal PIRTM factors (NEW)
  tierFactor:       number;   // T_i
  kernelConfidence: number;   // c_i
}

export interface ContributionWeight {
  orgId:   string;
  weight:  number;    // w_i^final = T_i · c_i · r_i
  factors: ContributionWeightFactors;

  // ADR-028: provenance (NEW)
  inferenceMethod:        InferenceMethod;
  kernelConfidenceMethod: KernelConfidenceMethod;
}
\end{lstlisting}

Existing callers that only rely on \texttt{weight} and the original factor
fields remain valid; the new fields are additive and defaulted to safe
values for Tier--1 cumulant paths when omitted.\,[file:12]

\subsection{Reputation Engine Computation}

The extended \texttt{calculateContributionWeight} function now receives
explicit inference provenance and kernel error information, but defaults
preserve legacy behavior:

\begin{lstlisting}[language=TypeScript,caption={Three--factor weight computation in \texttt{ReputationEngine}.}]
async calculateContributionWeight(
  orgId:              string,
  inferenceMethod:    InferenceMethod = 'cumulant-primary',
  kernelErrorRatio?:  number,       // |ε| / |Ξ|
  hasTier1InEnsemble: boolean = false,
): Promise<ContributionWeight> {
  const reputation = await this.reputationStore.getReputation(orgId);
  if (!reputation) {
    throw new Error(`No reputation found for organization ${orgId}`);
  }

  // Existing reputation product r_i
  const baseReputation   = reputation.reputationScore;
  const stakeMultiplier  = this.calculateStakeMultiplier(reputation.stakePledge);
  const consistencyBonus = this.calculateConsistencyBonus(reputation.consistencyScore);
  const r_i              = baseReputation * (1 + stakeMultiplier) * (1 + consistencyBonus);

  // Tier factor T_i
  const tierFactor: number = TIER_FACTORS[inferenceMethod];

  // Kernel confidence c_i
  let kernelConfidence:       number;
  let kernelConfidenceMethod: KernelConfidenceMethod;

  if (
    inferenceMethod === 'cumulant-primary' ||
    inferenceMethod === 'pseudoinverse-secondary'
  ) {
    if (kernelErrorRatio === undefined) {
      kernelConfidence       = 1.0;
      kernelConfidenceMethod = 'direct';
    } else {
      kernelConfidence       = Math.max(0, 1 - kernelErrorRatio);
      kernelConfidenceMethod = 'direct';
    }
  } else {
    // Tier 3: depend on presence of Tier-1 anchor
    if (hasTier1InEnsemble && kernelErrorRatio !== undefined) {
      kernelConfidence       = Math.max(0, 1 - kernelErrorRatio);
      kernelConfidenceMethod = 'cross-validated';
    } else {
      kernelConfidence       = TIER_FACTORS['moorpenrose-tertiary']; // floor β
      kernelConfidenceMethod = 'conservative-floor';
    }
  }

  const weight = tierFactor * kernelConfidence * r_i;

  return {
    orgId,
    weight,
    factors: {
      baseReputation,
      stakeMultiplier,
      consistencyBonus,
      totalMultiplier: r_i,
      tierFactor,
      kernelConfidence,
    },
    inferenceMethod,
    kernelConfidenceMethod,
  };
}
\end{lstlisting}

This function encodes the ADR--028 constitutional principle: Tier--3 models
cannot claim high confidence unless they are anchored against a Tier--1
cumulant model in the same ensemble.\,[file:12]

\section{L0 Guard in Byzantine Filtering}

\subsection{Invariant: Tier--3 Cross--Validation Requires Tier--1 Anchor}

To enforce ADR--028 at runtime, a guard is introduced into
\texttt{ByzantineFilter.prepareWeightedContributions}. The invariant is:

\begin{quote}
  If a Tier--3 model claims cross--validated kernel confidence
  (\texttt{kernelConfidenceMethod = 'cross-validated'}), the ensemble must
  contain at least one Tier--1 (\texttt{'cumulant-primary'}) contributor; otherwise,
  downgrade the confidence to the conservative floor and relabel the method
  as \texttt{'conservative-floor'}.\,[file:12]
\end{quote}

\subsection{Guard Implementation}

\begin{lstlisting}[language=TypeScript,caption={L0 guard in \texttt{prepareWeightedContributions}.}]
private prepareWeightedContributions(
  contributions: RawContribution[],
  weights:       Map<string, ContributionWeight>,
): WeightedContribution[] {
  return contributions.map(contrib => {
    const weightData = weights.get(contrib.orgIdHash);
    const weight     = weightData?.weight ?? 0.5;
    const factors    = weightData?.factors ?? {
      baseReputation:  0.5,
      stakeMultiplier: 0,
      consistencyBonus: 0,
      totalMultiplier: 1.0,
      tierFactor:       TIER_FACTORS['moorpenrose-tertiary'],
      kernelConfidence: TIER_FACTORS['moorpenrose-tertiary'],
    };

    // L0 INVARIANT: Tier-3 cross-validation requires Tier-1 anchor.
    if (
      weightData?.inferenceMethod        === 'moorpenrose-tertiary' &&
      weightData?.kernelConfidenceMethod === 'cross-validated'
    ) {
      const ensembleHasTier1 = Array.from(weights.values()).some(
        w => w.inferenceMethod === 'cumulant-primary'
      );
      if (!ensembleHasTier1) {
        factors.kernelConfidence = TIER_FACTORS['moorpenrose-tertiary'];
        (weightData as any).kernelConfidenceMethod = 'conservative-floor';
        console.warn(
          '[L0-ADR-028] Tier-3 cross-validated confidence downgraded: ' +
          `no Tier-1 anchor present in ensemble for org ` +
          `${contrib.orgIdHash.substring(0, 8)}…`
        );
      }
    }

    return {
      orgIdHash:     contrib.orgIdHash,
      fpRate:        contrib.fpRate,
      weight,
      eventCount:    contrib.eventCount,
      zScore:        0,
      weightFactors: factors,
    };
  });
}
\end{lstlisting}

This guard sits at the boundary of the trust plane and the compute plane, and
is explicitly treated as an L0 invariant: no higher--layer aggregation or
calibration is allowed to override it.\,[file:6][file:12]

\section{Monoidal Ensemble Aggregation in Python}

\subsection{Pairwise Commutator Signals}

On the Python side, each contribution carries both its final weight and its
inferred operator \(\Xi_i\) encoded as a matrix. The pairwise emergent signal
\(E_{ij}\) is constructed as:

\begin{lstlisting}[language=Python3,caption={Pairwise emergent signal \(E_{ij}\).}]
from __future__ import annotations
from dataclasses import dataclass
from itertools import combinations
from functools import reduce
from typing import Any
import numpy as np

def compute_final_weight(contribution: dict) -> float:
    f = contribution["factors"]
    return f["tierFactor"] * f["kernelConfidence"] * f["totalMultiplier"]

def compute_E_ij(ci: dict, cj: dict) -> dict:
    """
    E_ij(t) = w_i^final · w_j^final · [Ξ_i, Ξ_j]
    """
    w_i = compute_final_weight(ci)
    w_j = compute_final_weight(cj)

    xi_i = np.array(ci["xi_operator"])  # Ξ_i
    xi_j = np.array(cj["xi_operator"])  # Ξ_j

    commutator = xi_i @ xi_j - xi_j @ xi_i
    comm_mag   = float(np.linalg.norm(commutator, ord="fro"))

    return {
      "signal":     commutator * w_i * w_j,
      "magnitude":  comm_mag * w_i * w_j,
      "weight":     w_i * w_j,
      "provenance": {ci["modelId"], cj["modelId"]},
      "inferenceMethodPair": (
          ci["inferenceMethod"],
          cj["inferenceMethod"],
      ),
    }
\end{lstlisting}

The key multiplicity--theoretic point is that the commutator remains in
\(B(H)\), so multiplication by weights and later addition do not require any
projection.\,[file:12]

\subsection{Monoidal Product and Ensemble Reduction}

The monoidal product over emergent signals is defined as simple addition in
the operator component, multiplication of weights, and set union of
provenance:

\begin{lstlisting}[language=Python3,caption={Monoidal product and ensemble aggregation.}]
def _identity_signal(operator_shape: tuple) -> dict:
    return {
      "signal":     np.zeros(operator_shape),
      "magnitude":  0.0,
      "weight":     1.0,
      "provenance": set(),
      "inferenceMethodPair": (),
    }

def monoidal_product(e_ij: dict, e_kl: dict) -> dict:
    """
    Associative combination of pairwise emergent signals.
    """
    return {
      "signal":     e_ij["signal"] + e_kl["signal"],
      "magnitude":  e_ij["magnitude"] + e_kl["magnitude"],
      "weight":     e_ij["weight"] * e_kl["weight"],
      "provenance": e_ij["provenance"] | e_kl["provenance"],
      "inferenceMethodPair": (),
    }

def aggregate_ensemble(contributions: list[dict]) -> dict:
    if len(contributions) < 2:
        raise ValueError("Ensemble requires at least 2 contributions.")

    pairs = list(combinations(contributions, 2))
    pairwise_signals = [compute_E_ij(ci, cj) for ci, cj in pairs]

    op_shape = pairwise_signals[0]["signal"].shape
    identity = _identity_signal(op_shape)

    ensemble_signal = reduce(monoidal_product, pairwise_signals, identity)

    return {
      **ensemble_signal,
      "pairCount":        len(pairs),
      "contributorCount": len(contributions),
      "attributedModels": list(ensemble_signal["provenance"]),
    }
\end{lstlisting}

The associativity of this reduction follows directly from the associativity of
addition in \(B(H)\); thus, the monoidal structure is given ``for free'' by
the Lie--algebraic setting.\,[file:12]

Downstream systems (e.g.\ \texttt{originRecord}, FP calibration dashboards)
can store the aggregated operator, its Frobenius norm, composite trust weight,
and the provenance set \texttt{attributedModels}, enabling governance--aware
interpretation of ensemble behavior.\,[file:13][file:15]

\section{Governance and Orchestration Implications}

\subsection{Central Tension}

The central tension exposed and resolved here is:

\begin{quote}
  \emph{Accuracy vs.\ governance in trust weighting:}
  how to admit heterogeneous models (tiered inference provenance) into a
  single ensemble, while preserving constitutional guarantees about how
  confidence is earned, anchored, and enforced.
\end{quote}

On one side, the multiplicity--theoretic Lie algebra structure encourages a
rich operator space with many interacting signals; on the other, the
governance plane demands that not all operators are trusted equally or in the
same way.\,[file:6][file:13]

\subsection{Levers, Metrics, and Horizons}

Consistent with the Phase Mirror methodology, the developments above can be
cast as levers:

\begin{itemize}[leftmargin=*]
  \item \textbf{Lever 1 (Eng):} Enforce ADR--028's three--factor weight model
  in the TypeScript trust module.
  \emph{Metric:} all calibration paths use \(w_i^{\text{final}} =
  T_i c_i r_i\); no legacy code paths compute weights differently.
  \emph{Horizon:} 7 days.\,[file:12]

  \item \textbf{Lever 2 (Eng):} Enforce L0 guard for Tier--3 cross--validation.
  \emph{Metric:} any Tier--3 cross--validated confidence observed without a
  Tier--1 anchor triggers downgrade and a logged L0 warning.
  \emph{Horizon:} 7 days.\,[file:12][file:6]

  \item \textbf{Lever 3 (Lead Multiplicity Theorist):} Ratify ADR--028 and
  \texttt{SPEC-TRUST.md} as the constitutional reference for all weighting
  decisions.
  \emph{Metric:} ADR--028 cited by any new trust--related ADR or rule.
  \emph{Horizon:} 7--21 days.\,[file:6][file:12]

  \item \textbf{Lever 4 (Eng):} Wire Python monoidal reduction into the FP
  calibration pipeline (Plane B) with provenance reporting into the trust
  dashboards (Plane C).
  \emph{Metric:} ensemble calibration outputs include
  \texttt{attributedModels}, emergent--signal magnitude, and confidence
  metrics.
  \emph{Horizon:} 21 days.\,[file:12][file:15]
\end{itemize}

\subsection{Suggested Further Sections}

If you extend this report, natural additions include:

\begin{itemize}[leftmargin=*]
  \item A formal cold--start analysis of reputation convergence and the
  probation tier, connecting the three--factor weight model to CBFT/credit
  BFT literature.\,[file:6]
  \item A more detailed multiplicity--theoretic account of how prime--indexed
  operators \(U_p^{(r)}\) encode nonlinear emergent structure across scales.
  \item A specification section for \texttt{originRecord} schema evolution to
  store Lie--algebraic ensemble signals alongside scalar calibration metrics.
\end{itemize}

\appendix
\section*{Mathematical Appendix}
\addcontentsline{toc}{section}{Mathematical Appendix}

\section{Operator Algebra and Lie Structure}

\subsection{Boundedness and Algebra Structure of \texorpdfstring{$B(H)$}{B(H)}}

Let
\[
  H = \ell^2(\mathbb{P}) \otimes L^2(\mathbb{R}) \otimes \mathbb{C}^d
\]
denote the multiplicity Hilbert space used by PIRTM.\footnote{See trust module
and PIRTM integration summaries for the explicit construction of the Hilbert
components and their roles in the trust architecture.\,[file:12]} Denote by
\(B(H)\) the set of all bounded linear operators on \(H\).

\begin{lemma}
\(B(H)\) is a complex associative algebra with identity under composition and
a complex vector space under pointwise addition and scalar multiplication.
\end{lemma}

\begin{proof}
For any Hilbert space \(H\), the set \(B(H)\) of bounded linear operators is
closed under composition: if \(X,Y \in B(H)\) then \(XY\) is bounded and linear,
with
\[
  \|XY\| \le \|X\|\ \|Y\|.
\]
Associativity of composition follows from associativity of function
composition:
\((XY)Z = X(YZ)\) for all \(X,Y,Z \in B(H)\).

Moreover, \(B(H)\) is closed under pointwise addition and scalar
multiplication: if \(X,Y \in B(H)\) and \(\alpha,\beta \in \mathbb{C}\), then
\(\alpha X + \beta Y\) is also a bounded linear operator with
\(\|\alpha X + \beta Y\| \le |\alpha|\|X\| + |\beta|\|Y\|\).

Thus \(B(H)\) is a complex vector space and an associative algebra with
identity, given by the identity operator \(I_H\).
\end{proof}

In particular, each PIRTM operator \(U_p^{(r)}\) and any bundled operator
\(\Xi_i\) appearing in the trust pipeline is an element of \(B(H)\).\,[file:12]

\subsection{Commutator Defines a Lie Algebra Structure}

Define the commutator
\[
  [X,Y] := XY - YX
\]
for \(X,Y \in B(H)\).

\begin{proposition}
The pair \((B(H),[\cdot,\cdot])\) is a Lie algebra over \(\mathbb{C}\).
\end{proposition}

\begin{proof}
We verify the three Lie algebra axioms:

\textbf{(1) Bilinearity.}
For any \(\alpha,\beta \in \mathbb{C}\) and \(X_1,X_2,Y \in B(H)\),
\begin{align*}
  [\alpha X_1 + \beta X_2, Y]
  &= (\alpha X_1 + \beta X_2)Y - Y(\alpha X_1 + \beta X_2) \\
  &= \alpha(X_1Y - YX_1) + \beta(X_2Y - YX_2) \\
  &= \alpha[X_1,Y] + \beta[X_2,Y].
\end{align*}
A symmetric calculation shows linearity in the second argument.

\textbf{(2) Antisymmetry.}
For all \(X,Y \in B(H)\),
\[
  [X,Y] = XY - YX = - (YX - XY) = -[Y,X].
\]

\textbf{(3) Jacobi identity.}
For all \(X,Y,Z \in B(H)\), the identity
\[
  [X,[Y,Z]] + [Y,[Z,X]] + [Z,[X,Y]] = 0
\]
holds. This is a standard consequence of associativity and can be checked by
expanding each term:
\begin{align*}
[X,[Y,Z]]
&= X(YZ - ZY) - (YZ - ZY)X \\
&= XYZ - XZY - YZX + ZYX,\\
[Y,[Z,X]]
&= Y(ZX - XZ) - (ZX - XZ)Y \\
&= YZX - YXZ - ZXY + XZY,\\
[Z,[X,Y]]
&= Z(XY - YX) - (XY - YX)Z \\
&= ZXY - ZYX - XYZ + YXZ.
\end{align*}
Summing these three expansions, all terms cancel in pairs, so the sum is
zero.
\end{proof}

Consequently, any family of operators (e.g.\ the PIRTM atlas and its
finite--linear combinations) generates a Lie subalgebra of \(B(H)\) under the
commutator bracket.\,[file:12]

\subsection{Emergent Signals Lie in the Same Operator Space}

For each contributing model \(i\), consider an inferred operator
\(\Xi_i \in B(H)\) and a scalar weight \(w_i^{\mathrm{final}} \in \mathbb{R}\).
Define the pairwise emergent signal
\[
  \mathcal{E}_{ij} := w_i^{\mathrm{final}} w_j^{\mathrm{final}} [\Xi_i,\Xi_j].
\]

\begin{lemma}
For all \(i,j\), we have \(\mathcal{E}_{ij} \in B(H)\).
\end{lemma}

\begin{proof}
Since \(\Xi_i, \Xi_j \in B(H)\), we know \(\Xi_i\Xi_j, \Xi_j\Xi_i \in B(H)\).
Therefore \([\Xi_i,\Xi_j] = \Xi_i\Xi_j - \Xi_j\Xi_i \in B(H)\) by closure of
\(B(H)\) under addition and scalar multiplication.

Moreover, \(w_i^{\mathrm{final}} w_j^{\mathrm{final}}\) is a scalar, so
\(w_i^{\mathrm{final}} w_j^{\mathrm{final}} [\Xi_i,\Xi_j]\) is still a bounded
operator. Thus \(\mathcal{E}_{ij} \in B(H)\) for all \(i,j\).
\end{proof}

\begin{corollary}
Any finite linear combination of emergent signals
\(\sum_{(i,j)} \alpha_{ij}\,\mathcal{E}_{ij}\) with \(\alpha_{ij} \in \mathbb{C}\)
belongs to \(B(H)\).
\end{corollary}

\begin{proof}
Each \(\mathcal{E}_{ij} \in B(H)\), and \(B(H)\) is a vector space, so finite
linear combinations remain in \(B(H)\).
\end{proof}

This justifies the design decision in the ensemble aggregator: additive
reduction over pairwise commutators remains in the same ambient operator
space, with no need for projection back onto a restricted atlas.\,[file:12]

\section{Operator Norm Bounds for Commutators}

\subsection{Submultiplicative Bounds}

For \(X,Y \in B(H)\), recall the operator norm
\[
  \|X\| := \sup_{\|v\|_H = 1} \|Xv\|_H.
\]

\begin{lemma}[Basic commutator bound]
\label{lem:basic-commutator-bound}
For all \(X,Y \in B(H)\),
\[
  \|[X,Y]\| \le 2 \|X\|\,\|Y\|.
\]
\end{lemma}

\begin{proof}
Using the triangle inequality and submultiplicativity of the operator norm:
\begin{align*}
  \|[X,Y]\|
  &= \|XY - YX\| \\
  &\le \|XY\| + \|YX\| \\
  &\le \|X\|\,\|Y\| + \|Y\|\,\|X\| \\
  &= 2 \|X\|\,\|Y\|.
\end{align*}
\end{proof}

In many PIRTM/Phase Mirror implementations, the emergent signal is represented
by a finite matrix acting on a truncated basis of \(H\), and the Frobenius
norm is used as a practical surrogate for the operator norm.\,[file:12] In
finite dimensions, we have \(\|M\|_2 \le \|M\|_F\), so the bound above can be
translated into Frobenius--norm terms if needed.

\subsection{Weighted Commutator Bounds}

The emergent signal is weighted by the final trust weights. For each pair
\((i,j)\), define
\[
  \mathcal{E}_{ij} = w_i^{\mathrm{final}} w_j^{\mathrm{final}} [\Xi_i,\Xi_j].
\]

\begin{proposition}[Norm bound for weighted emergent signal]
\label{prop:weighted-bound}
Let \(M_i := \|\Xi_i\|\) and \(W_i := |w_i^{\mathrm{final}}|\). Then
\[
  \|\mathcal{E}_{ij}\| \le 2\,W_i W_j\,M_i M_j.
\]
\end{proposition}

\begin{proof}
By Lemma~\ref{lem:basic-commutator-bound},
\(\|[\Xi_i,\Xi_j]\| \le 2\,\|\Xi_i\|\|\Xi_j\| = 2 M_i M_j\). Multiplying by the
scalar \(w_i^{\mathrm{final}} w_j^{\mathrm{final}}\) yields
\[
  \|\mathcal{E}_{ij}\|
  = |w_i^{\mathrm{final}} w_j^{\mathrm{final}}|\,\|[\Xi_i,\Xi_j]\|
  \le (W_i W_j) \cdot (2 M_i M_j)
  = 2\,W_i W_j\,M_i M_j.
\]
\end{proof}

This bound can be used to enforce or diagnose L0--style bounds on the
magnitude of emergent signals relative to tier, kernel confidence, and
reputation.\,[file:6][file:12]

\section{Monoidal Ensemble Reduction: Associativity and Bounds}

\subsection{Associativity of the Ensemble Product}

The ensemble aggregation defines a monoidal product on emergent signals
\(\{E_{ij}\}\). In the Python implementation, an individual emergent signal is
represented as a record
\[
  E = (\Sigma, \mu, w, P)
\]
where:
\begin{itemize}[leftmargin=*]
  \item \(\Sigma\) is the operator--valued signal (matrix representation of
  \(B(H)\));
  \item \(\mu\) is a scalar magnitude (e.g.\ Frobenius norm of \(\Sigma\));
  \item \(w\) is a scalar weight (product of pairwise weights);
  \item \(P\) is a provenance set of model identifiers.\,[file:12][file:13]
\end{itemize}

The monoidal product \(\otimes\) is given by:
\[
  E \otimes E' =
  (\Sigma + \Sigma',\ \mu + \mu',\ w w',\ P \cup P').
\]

\begin{proposition}
The product \(\otimes\) is associative.
\end{proposition}

\begin{proof}
Let \(E_1 = (\Sigma_1,\mu_1,w_1,P_1)\),
\(E_2 = (\Sigma_2,\mu_2,w_2,P_2)\),
\(E_3 = (\Sigma_3,\mu_3,w_3,P_3)\).

Compute \((E_1 \otimes E_2) \otimes E_3\):
\begin{align*}
  E_1 \otimes E_2
  &= (\Sigma_1 + \Sigma_2,\ \mu_1 + \mu_2,\ w_1 w_2,\ P_1 \cup P_2),\\
  (E_1 \otimes E_2) \otimes E_3
  &= ((\Sigma_1 + \Sigma_2) + \Sigma_3,\ (\mu_1 + \mu_2) + \mu_3,\\
  &\quad (w_1 w_2) w_3,\ (P_1 \cup P_2) \cup P_3).
\end{align*}

Similarly, \(E_1 \otimes (E_2 \otimes E_3)\) equals
\begin{align*}
  E_2 \otimes E_3
  &= (\Sigma_2 + \Sigma_3,\ \mu_2 + \mu_3,\ w_2 w_3,\ P_2 \cup P_3),\\
  E_1 \otimes (E_2 \otimes E_3)
  &= (\Sigma_1 + (\Sigma_2 + \Sigma_3),\ \mu_1 + (\mu_2 + \mu_3),\\
  &\quad w_1 (w_2 w_3),\ P_1 \cup (P_2 \cup P_3)).
\end{align*}

Using associativity of addition in the vector space of operators,
associativity of addition in \(\mathbb{R}\), associativity of scalar
multiplication, and associativity of set union, we have:
\begin{align*}
(\Sigma_1 + \Sigma_2) + \Sigma_3 &= \Sigma_1 + (\Sigma_2 + \Sigma_3),\\
(\mu_1 + \mu_2) + \mu_3 &= \mu_1 + (\mu_2 + \mu_3),\\
(w_1 w_2) w_3 &= w_1 (w_2 w_3),\\
(P_1 \cup P_2) \cup P_3 &= P_1 \cup (P_2 \cup P_3).
\end{align*}
Hence \((E_1 \otimes E_2) \otimes E_3 = E_1 \otimes (E_2 \otimes E_3)\).
\end{proof}

The monoidal unit is the zero signal
\[
  E_{\mathrm{id}} := (0,\ 0,\ 1,\ \varnothing),
\]
where \(0\) is the zero operator and the scalar zero. One easily checks that
\(E \otimes E_{\mathrm{id}} = E_{\mathrm{id}} \otimes E = E\).

\subsection{Norm Bounds for Ensemble Signal}

Let \(\{E_{ij}\}\) be the collection of pairwise emergent signals produced by
an ensemble of \(n\) models:
\[
  E_{ij} = (\Sigma_{ij},\mu_{ij},w_{ij},P_{ij}), \quad 1 \le i < j \le n.
\]
Their monoidal reduction is
\[
  E_{\mathrm{ensemble}} =
  \bigotimes_{1 \le i < j \le n} E_{ij}
  = \left(
      \sum_{i<j} \Sigma_{ij},\
      \sum_{i<j} \mu_{ij},\
      \prod_{i<j} w_{ij},\
      \bigcup_{i<j} P_{ij}
    \right).
\]

\begin{proposition}[Ensemble operator norm bound]
\label{prop:ensemble-bound}
Let \(\|\cdot\|\) be an operator norm on \(B(H)\), and suppose
\(\|\Sigma_{ij}\| \le K_{ij}\) for all \(i<j\). Then
\[
  \left\|\sum_{i<j} \Sigma_{ij}\right\|
  \le \sum_{i<j} K_{ij}.
\]
In particular, if each \(\Sigma_{ij} = w_i^{\mathrm{final}} w_j^{\mathrm{final}} [\Xi_i,\Xi_j]\)
satisfies the bound from Proposition~\ref{prop:weighted-bound}, then
\[
  \bigl\|\Sigma_{\mathrm{ensemble}}\bigr\|
  \le 2\sum_{i<j} W_i W_j M_i M_j.
\]
\end{proposition}

\begin{proof}
The first inequality is just the triangle inequality applied repeatedly:
\begin{align*}
  \left\|\sum_{i<j} \Sigma_{ij}\right\|
  &\le \sum_{i<j} \|\Sigma_{ij}\|
  \le \sum_{i<j} K_{ij}.
\end{align*}

For the second claim, take
\(K_{ij} := 2 W_i W_j M_i M_j\) as in
Proposition~\ref{prop:weighted-bound}. Then
\(\|\Sigma_{ij}\| \le K_{ij}\), and substituting into the first inequality
gives the desired bound.
\end{proof}

This provides an explicit link between per--model operator norms, trust
weights, and the maximal magnitude of the ensemble emergent signal. It can be
used to define L0 or L1 thresholds on acceptable ensemble behavior (e.g.\ by
capping \(\bigl\|\Sigma_{\mathrm{ensemble}}\bigr\|\) relative to tier
composition or kernel--confidence profiles).\,[file:6][file:12]

\section{Weight Model Constraints and Bounds}

\subsection{Three--Factor Weight Model}

ADR--028 defines
\[
  w_i^{\mathrm{final}} = T_i \cdot c_i \cdot r_i,
\]
where \(T_i\) is a tier factor,
\(c_i \in [0,1]\) is a kernel--confidence factor,
and \(r_i \ge 0\) is a reputation--based multiplier.\,[file:12]

Assume the following constraints, consistent with the implementation:
\begin{itemize}[leftmargin=*]
  \item Tier factors are bounded: \(0 \le T_i \le 1\).
  \item Kernel confidence is bounded: \(0 \le c_i \le 1\).
  \item Reputation multipliers satisfy \(0 \le r_i \le R_{\max}\), where
  \(R_{\max}\) depends on the maximum reputation, stake multiplier, and
  consistency bonus allowed by policy.\,[file:12][file:15]
\end{itemize}

\begin{lemma}[Per--model weight bound]
\label{lem:wi-bound}
Under the above constraints,
\[
  0 \le w_i^{\mathrm{final}} \le R_{\max}.
\]
\end{lemma}

\begin{proof}
Since \(T_i, c_i, r_i \ge 0\), clearly \(w_i^{\mathrm{final}} \ge 0\).
Moreover, \(T_i \le 1\) and \(c_i \le 1\), so
\[
  w_i^{\mathrm{final}} = T_i c_i r_i \le 1 \cdot 1 \cdot r_i \le R_{\max}.
\]
\end{proof}

\subsection{Pairwise and Ensemble Weight Bounds}

The pairwise product weight is \(w_{ij} := w_i^{\mathrm{final}} w_j^{\mathrm{final}}\).

\begin{lemma}[Pairwise weight bound]
For all \(i \ne j\),
\[
  0 \le w_{ij} \le R_{\max}^2.
\]
\end{lemma}

\begin{proof}
From Lemma~\ref{lem:wi-bound}, \(0 \le w_i^{\mathrm{final}} \le R_{\max}\) and
\(0 \le w_j^{\mathrm{final}} \le R_{\max}\). Hence
\[
  0 \le w_{ij} = w_i^{\mathrm{final}} w_j^{\mathrm{final}} \le R_{\max}^2.
\]
\end{proof}

For an ensemble of \(n\) models, the aggregate weight in the monoidal
reduction is
\[
  w_{\mathrm{ensemble}} := \prod_{1 \le i < j \le n} w_{ij}.
\]

\begin{proposition}[Ensemble weight bound]
Assuming all \(w_{ij} \le R_{\max}^2\) and \(w_{ij} \ge 0\), we have
\[
  0 \le w_{\mathrm{ensemble}} \le R_{\max}^{n(n-1)}.
\]
\end{proposition}

\begin{proof}
There are \(\binom{n}{2} = \frac{n(n-1)}{2}\) distinct pairs \((i,j)\) with
\(i<j\). Thus
\[
  w_{\mathrm{ensemble}}
  = \prod_{i<j} w_{ij}
  \le \prod_{i<j} R_{\max}^2
  = R_{\max}^{2 \cdot \binom{n}{2}}
  = R_{\max}^{n(n-1)}.
\]
Nonnegativity follows from nonnegativity of each factor.
\end{proof}

In practice, the implementation may normalize or re--parameterize the
ensemble weight (e.g.\ via logarithms or clipping) to avoid numerical
overflow or to reflect governance--driven decay semantics, but the bound
above provides a worst--case envelope anchored to the underlying trust
factors.\,[file:12][file:15]

\section{L0 Guard Correctness}

\subsection{Constraint: Tier--3 Cross--Validation Requires Tier--1 Anchor}

The L0 constraint enforced in the TypeScript code can be stated as:

\begin{quote}
If a contribution with \(\texttt{inferenceMethod} = \texttt{'moorpenrose-tertiary'}\)
claims \(\texttt{kernelConfidenceMethod} = \texttt{'cross-validated'}\), then
there must exist at least one contribution in the same ensemble with
\(\texttt{'cumulant-primary'}\) inference method. Otherwise, the
kernel--confidence must be reset to the Tier--3 floor and the method relabeled
\texttt{'conservative-floor'}.\,[file:12][file:6]
\end{quote}

\begin{proposition}[L0 guard enforces the constraint]
Let \(W\) be the multiset of contribution weight records for a given ensemble.
The L0 guard implemented in \texttt{prepareWeightedContributions} ensures that
in the resulting set of weights \(W'\), the above condition holds for all
Tier--3 contributions.
\end{proposition}

\begin{proof}
Consider any contribution \(c\) with weight record \(w_c \in W\) such that
\(\texttt{inferenceMethod}(w_c) = \texttt{'moorpenrose-tertiary'}\) and
\(\texttt{kernelConfidenceMethod}(w_c) = \texttt{'cross-validated'}\).

The guard computes
\[
  \texttt{ensembleHasTier1}
  := \exists w \in W\ \text{such that } \texttt{inferenceMethod}(w)
      = \texttt{'cumulant-primary'}.
\]
If \(\texttt{ensembleHasTier1} = \mathrm{false}\), it updates \(w_c\) (and the
local factors) so that:
\[
  \texttt{kernelConfidence}(w_c) \gets T_{\mathrm{Tier3}},
  \quad
  \texttt{kernelConfidenceMethod}(w_c) \gets \texttt{'conservative-floor'}.
\]
Thus in the resulting set \(W'\) of weight records after the guard:

\begin{itemize}[leftmargin=*]
  \item Either there exists at least one Tier--1 contribution in the ensemble,
  in which case cross--validated Tier--3 confidence is permitted.
  \item Or every contribution with Tier--3 inference method has
  \(\texttt{kernelConfidenceMethod} \ne \texttt{'cross-validated'}\), because
  such entries were downgraded to \texttt{'conservative-floor'}.
\end{itemize}

Therefore the constraint holds in \(W'\).
\end{proof}

This is a formal statement that the implementation matches the constitutional
intent encoded in ADR--028, and that no higher--layer mechanism can override
the L0 restriction without directly editing this guard.\,[file:6][file:12]

\section{Summary of Key Inequalities}

For quick reference, the main operator and weight bounds established in this
appendix are:

\begin{itemize}[leftmargin=*]
  \item Commutator norm bound:
  \[
    \|[X,Y]\| \le 2 \|X\|\|Y\|.
  \]
  \item Weighted emergent signal:
  \[
    \|\mathcal{E}_{ij}\|
    = \|w_i^{\mathrm{final}} w_j^{\mathrm{final}} [\Xi_i,\Xi_j]\|
    \le 2\,W_i W_j\,M_i M_j.
  \]
  \item Ensemble signal:
  \[
    \left\|\sum_{i<j} \Sigma_{ij}\right\|
    \le \sum_{i<j} \|\Sigma_{ij}\|
    \le 2\sum_{i<j} W_i W_j M_i M_j.
  \]
  \item Per--model weight:
  \[
    0 \le w_i^{\mathrm{final}} \le R_{\max}.
  \]
  \item Pairwise and ensemble weights:
  \[
    0 \le w_{ij} \le R_{\max}^2, \qquad
    0 \le w_{\mathrm{ensemble}} \le R_{\max}^{n(n-1)}.
  \]
\end{itemize}

These inequalities, together with the associativity of the monoidal product
and the L0 guard, provide the formal backbone for the trust--weighted
Lie--algebraic aggregation used in the Phase Mirror / PIRTM integration.\,[file:12][file:6][file:15]


\nocite{*}
\bibliographystyle{unsrt}  
\bibliography{references}
\end{document}